\subsection*{从微观移动到宏观边界条件}

导体与绝缘体的区别，不在于材料中“有没有电荷”，而在于是否存在能在所研究时间尺度内作宏观
迁移的载流子。金属中的自由电子只要受到很小的切向电场，就会产生定向运动；因此静电平衡不是
外电场突然消失，而是载流子经过一个暂态重排后，表面电荷产生的附加电场恰好把导体内部的总场
抵消。绝缘体中的束缚电荷也会发生微小位移或取向极化，但通常不能像金属自由电子那样在整块材料
中重新分配。

由此可把导体静电平衡的结论串成一条逻辑链：导体材料内部 $\vect E=0$；于是导体内部电势没有
空间变化；连通导体连同表面成为一个等势体；表面切向场为零，外侧场只能沿法线；再用跨过表面
的薄柱高斯面，才得到 $E_n^{\rm(out)}=\sigma/\varepsilon_0$。这些结论必须按此顺序理解，不能
把“等势”误解成“表面没有电荷”，也不能把局部法向场直接套成球对称场。

图 31.1(a) 中，两块导体各自等势，但二者电势可以不同。实线从高电势导体出发并沿法线离开表面，
进入低电势导体时仍与表面正交；虚线表示等势线，处处与电场线正交。两导体之间的电场是一个
整体边值问题，不能把每块导体都当成互不影响的孤立带电球。表面某一点电荷密度的正负，应由该点
外侧场的法向方向判断，而不能只由整块导体的净电荷符号判断。

\subsection*{感应起电的全过程}

以带正电的外物体靠近中性金属球为例。最初，电子向近端移动，近端带负电、远端带等量正电，
金属球总电荷仍为零。保持外电荷位置不变并接地后，大地与球构成一个巨大的连通导体；在外正
电荷吸引下，电子从大地流入球体，原先远端的正电荷被中和。此时球的电势被固定为大地电势，
并不要求球的净电荷等于零。

随后必须先断开地线，使已经流入球内的电子失去返回大地的通道；最后再移去外电荷。外场撤除后，
球面上的电子在相互排斥下重新分布。若球是孤立球面且周围没有其他物体，最终分布具有球对称性；
若导体形状不规则或附近仍有其他导体，最终电荷仍会不均匀。“均匀分布”是球对称边界条件的结果，
不是感应起电对任意导体都成立的结论。若把最后两步颠倒，在仍然接地时先撤去外电荷，多余电子会
流回大地，球通常重新变为中性。

\subsection*{空腔与静电屏蔽的边界}

闭合导体空腔内没有自由电荷时，腔内满足 Laplace 方程，而整个腔壁属于同一连通导体、具有同一
电势。常量电势显然满足方程与边界条件；唯一性定理保证它就是唯一解，所以腔内电场为零。这个
结论并不需要腔体是球形，也不需要外界电场具有任何对称性。

若腔内放入电荷 $q$，腔内不再是零场。在导体材料中作包围空腔的高斯面，由于该面上场为零，
腔内壁的总感应电荷必为 $-q$。若孤立导体原有净电荷 $Q$，则外表面的总电荷为 $Q+q$；若导体
接地，外表面总电荷由零电势边界和外界环境决定。这里“屏蔽”表示内外边值问题受导体等势边界
约束，并不表示金属壁把电场线当作实体挡住。

真实 Faraday 笼只在静电或足够低频的近似下可直接用上述结论。时变场会在导体中激发电流，
屏蔽效果取决于电导率、壁厚、趋肤深度、网孔尺寸和电磁波频率。图中的放电通道终止在笼外表面，
说明电荷主要沿导体外部重新分布；它不能推出任意金属网对任意频率都能实现完全屏蔽。

\subsection*{从 Laplace 方程到离散电势}

一维无电荷区域满足 $\dd^2\varphi/\dd x^2=0$。把相邻网格点间距记为 $a$，在 $x_n$ 两侧作
Taylor 展开：
\begin{align*}
 \varphi_{n+1}&=\varphi_n+a\varphi'_n+\frac{a^2}{2}\varphi''_n+O(a^3),\\
 \varphi_{n-1}&=\varphi_n-a\varphi'_n+\frac{a^2}{2}\varphi''_n+O(a^3).
\end{align*}
两式相加，奇次项相消，于是
\begin{equation}
 \varphi''(x_n)=\frac{\varphi_{n+1}-2\varphi_n+\varphi_{n-1}}{a^2}+O(a^2).
\end{equation}
在二维等间距方格上，$\nabla^2\varphi=0$ 相应变成
\begin{equation}
 \boxed{\varphi_{i,j}=\frac14\bigl(\varphi_{i+1,j}+\varphi_{i-1,j}
 +\varphi_{i,j+1}+\varphi_{i,j-1}\bigr)}.
\end{equation}
也就是说，无电荷区域内部一点的电势等于四个最近邻点电势的平均值。这个“平均值性质”解释了
为什么调和函数不能在内部产生孤立的最大值或最小值，也给出逐点迭代求复杂导体电势的数值方法。
图 31.4(b) 中标出的中心二阶差分正是连续 Laplace 方程与离散电路、数值网格之间的桥梁。

唯一性定理在此还提供了验算标准：候选电势不仅要让导体表面“看起来像等势面”，还必须在整个
求解区域满足正确的 Poisson/Laplace 方程，并满足所有边界以及无穷远条件。只检验其中一部分，
不能证明构造解就是物理解。

\subsection*{球面像电荷的几何关系}

接地球问题中，真实电荷 $Q$ 位于距球心 $L$ 处，像电荷放在同一半径上、距球心
$l=R^2/L$ 处。任取球面点 $P$，由反演几何可得
\begin{equation}
 \frac{PQ}{Pq'}=\frac{L}{R}.
\end{equation}
选择 $q'=-(R/L)Q$ 后，$P$ 点两项电势恰好抵消，因而整个球面而非某一个特殊点都保持零电势。
这也是图 31.3 中像电荷必须位于球内且靠近真实电荷一侧的原因。若图中把像电荷画在错误位置，
就会破坏上述恒定距离比，边界条件不再成立。

像电荷法只保证所选真实区域内的电势正确。球内的像电荷不是金属中真实存在的点电荷；真实感应
电荷连续分布在球面，其近 $Q$ 一侧密度较大。计算 $Q$ 所受力时，可直接用像电荷的库仑力，因为
此时只需感应场在 $Q$ 位置的瞬时值。计算装配能量时必须写成
$U=Q\varphi_{\rm ind}/2$，因为感应电荷是在 $Q$ 从无穷远移近的过程中逐渐建立的。用
$F=-\dd U/\dd L$ 应当能恢复同一个吸引力，这是一项重要自检。

接地球和孤立球是不同边值问题。接地球可以与大地交换电荷，只要求球面电势为零；孤立中性球
既要等势，又要保证表面总感应电荷为零，通常需在球心补一个像电荷。若孤立球本来带给定净电荷，
球心像电荷的大小还要相应改变。因此看到“导体球”时，必须先判断它是否接地、是否孤立以及净
电荷是否给定，再选择像电荷组。

\subsection*{等离子体频率与电离层传播}

电子云相对离子背景移动距离 $\xi$ 后，两端出现的等效面电荷产生回复场，在线性近似下得到
$\omega_p=\sqrt{ne^2/(m\varepsilon_0)}$。这一频率描述自由载流子集体重排的固有时间尺度：
密度越高，回复场越强；质量越大，响应越慢。电子质量远小于离子质量，所以通常先讨论电子
等离子体振荡。

对简化的无碰撞等离子体，横向电磁波的色散关系可写为
\begin{equation}
 \omega^2=\omega_p^2+c^2k^2.
\end{equation}
当 $\omega<\omega_p$ 时，传播常数成为虚数，波在等离子体中迅速衰减，宏观上表现为反射；当
$\omega>\omega_p$ 时，波可以传播。电离层中的电子密度随高度、昼夜和太阳活动变化，因此反射
并不发生在一张固定的“镜面”上。图 31.6(a) 用弯曲路径表示低频无线电波可返回地面、支持远距离
通信，同时以穿透路径提醒读者：更高频率的波可能越过电离层。实际传播还会受到碰撞、地磁场、
分层结构和入射角的影响，等离子体频率只是判断能否传播的第一尺度。

\subsection*{导体系统的电势与电荷关系}

有多块彼此分离的导体时，每块导体内部仍是等势体，但各导体电势一般不同。在线性静电范围内，
导体电荷和电势可写成
\[
 Q_i=\sum_j C_{ij}\varphi_j,
\]
其中系数只由几何形状、相对位置和介质决定。对角元为正，非对角元通常为负；改变一块导体电势
会通过感应改变其他导体的表面电荷分布。这一形式提醒我们：两导体的净电荷大小不能单独决定谁的
电势更高，导体尺寸和相互位置同样重要。

判断表面是否可能出现异号电荷时，可结合最高、最低电势原理。在无自由体电荷区域，调和函数的
极值只能出现在边界。电势最高的导体表面，外法向电场不能指向导体内部，因而其表面不会出现
负电荷区；电势最低者相应不会出现正电荷区。但一个净电荷为正、却不是最高电势的导体，面对更
高电势导体的一侧仍可能出现局部负电荷。局部符号与净电荷符号并非同一问题。

\subsection*{边值题的检查清单}

构造静电解时，应逐项核对以下内容：真实求解区域内的电荷源是否正确；每一块接地导体是否为零
电势；每一块孤立导体是否等势且总电荷满足题设；不同边界相交处是否同时满足条件；无穷远处的
电势或场是否具有正确衰减。像电荷的数量少并不意味着解更“物理”，只有全部条件同时满足才可
借唯一性定理确认。

对两个相互垂直的接地平面，真实电荷关于每个平面各产生一个异号像，再对两个平面同时反射产生
一个同号像。四个电荷组成的势在两平面上分别成对抵消。若漏掉双重反射的同号像，其中一个平面
可能仍为零势，另一个却不满足边界。半球凸起或球壳内电荷问题同理：应先明确真实区域，再把像
电荷全部放到真实区域之外，并逐个验证平面反射和球面反演关系。

图中位置关系可用无量纲比值复核。接地球外点电荷必须满足
\[
 \frac{l}{R}=\frac{R}{L}<1,
\]
所以像电荷必在球内；球壳内点电荷 $x<R$ 的反演像位于 $R^2/x>R$，必在球外。只要图把两者画
反，即使公式文字正确，读者也会得到相反的空间判断。

\subsection*{数值解与物理解的对应}

复杂导体形状往往没有简单解析像电荷。把求解区域划分为网格，固定边界节点电势，再反复用
“内部节点等于邻点平均值”更新，就可逐步逼近 Laplace 方程的解。迭代收敛后，电场由相邻节点
电势差近似得到，导体表面电荷密度再由法向场
$\sigma=\varepsilon_0E_n$ 求出。网格越细，几何边界的表示越准确，但计算量也越大；尖角附近
场变化剧烈，通常需要局部加密。

数值图同样要接受物理检查：等势线不得相交；电场线应与导体表面正交；无体电荷区域不能出现
孤立电势极值；由包围导体的高斯面计算出的通量应与总电荷一致。数值收敛只说明离散方程被满足，
若边界条件输入错误，仍会稳定收敛到错误答案。

\subsection*{场能、静电压强与机械平衡}

导体表面外侧电场为 $E=\sigma/\varepsilon_0$，但某个表面电荷元受到的不是这个包含自场的总场。
去除自场后，其他电荷在表面处提供平均场 $E_{\rm other}=E/2$，因此表面受到向外的静电压强
\[
 p=\frac{\sigma^2}{2\varepsilon_0}
 =\frac12\varepsilon_0E^2.
\]
这个结果与真空电场能量密度相同。带电薄球壳有向外膨胀趋势，电容器两极板则相互吸引；方向
应结合两侧场强和边界法线判断，不能只看 $\sigma^2$ 的正值。

若导体由电源维持固定电势，机械位移时电源会输送或吸收电荷，单看电场能变化不足以给出机械力，
还要计入电源做功。孤立导体保持电荷不变，接地导体保持电势不变，这两类约束看似只差一句话，
却可能使能量法中的控制变量完全不同。先写清外界保持的量，再对几何参数求导，是避免符号错误
的关键。
